\subsection{恒参或随参信道特性对信号传输的影响}
\subsubsection{恒参信道}
恒参信道是\textit{线性时不变系统}, 其传输特性可以表示为
\begin{equation}
    H(\omega)=|H(\omega)|e^{-j\varphi(\omega)}.
\end{equation}

\textbf{理想恒参信道}\quad $H(\omega)=Ke^{-j\omega t_d}$.
\begin{equation}
    \begin{cases}
        |H(\omega)|=K,              \\
        \varphi(\omega)=\omega t_d, \\
        \tau(\omega)=\frac{\mathrm{d}\varphi(\omega)}{\mathrm{d}\omega}=t_d. \qquad (\text{群时延特性})
    \end{cases}
\end{equation}

理想恒参信道的冲激响应如下
\begin{equation}
    h(t)=K\delta(t-t_d).
\end{equation}
则其输出为
\begin{equation}
    s_o(t)=Ks(t-t_d).
\end{equation}
这被称为\textbf{无失真传输}.

对于多理想恒参信道下复杂信号 $f(t)$ 的传输, 设两条信道传输衰减 $K$, 传输时延分别为 $\tau_1$ 和 $\tau_1+\tau$, 则
\begin{equation}
    f_o(t)=Kf(t-\tau_1)+Kf(t-\tau_1-\tau).
\end{equation}
其频谱为
\begin{equation}
    F_o(\omega)=KF(\omega)e^{-j\omega\tau_1}+KF(\omega)e^{-j\omega(\tau_1{\color{red} +}\tau)}.
\end{equation}
则此时的信道传输函数为
\begin{equation}
    H(\omega)=\frac{F_o(\omega)}{F(\omega)}=Ke^{-j\omega\tau_1}(1+e^{-j\omega\tau}).
\end{equation}
信道幅频特性
\begin{equation}
    |H(\omega)|=2K\left|\cos\frac{\omega\tau}{2}\right|.
\end{equation}

\textbf{频率选择性衰落}\quad 信道对信号不同频率部分有不同衰减.

减小频率选择性衰落影响的方法: 降低信号带宽
\begin{equation}
    \frac{1}{3}\Delta f\leq B\geq\frac{1}{5}\Delta f.
\end{equation}

\subsubsection{随参信道}
\textbf{特点}
\begin{itemize}
    \item 衰减和时延随时间变化;
    \item 多径传播.
\end{itemize}

\textbf{多径效应}

对于发送信号 $s(t)=A\cos\omega_ct$, 经过 $n$ 条路径传播, 结果为
\begin{equation}
    \begin{aligned}
        r(t) & =\sum_{i=1}^{n}A_i(t)\cos\omega_c(t-\tau_i(t))     \\
             & =\sum_{i=1}^{n}A_i(t)\cos(\omega_ct+\varphi_i(t)).
    \end{aligned}
\end{equation}

展开之后, 得到
\begin{equation}
    \begin{aligned}
        r(t) & =\sum_{i=1}^{n}A_i(t)\cos\varphi_i\cos\omega_ct-\sum_{i=1}^{n}A_i(t)\sin\varphi_i\sin\omega_ct \\
             & =X(t)\cos\omega_ct-Y(t)\sin\omega_ct.
    \end{aligned}
\end{equation}
显然, 这满足\textit{窄带随机过程}的形式.

多径传播使信号产生\textbf{瑞利衰落}, 引起\textbf{频率弥散}.
